\documentclass[manuscript]{acmart}

\setcopyright{none}
\settopmatter{printacmref=false}
\raggedbottom

\usepackage{booktabs}
\usepackage{multirow}
\usepackage{enumitem}
\usepackage{xspace}
\usepackage{graphicx}
\usepackage{hyperref}
\usepackage{microtype}
\usepackage{xcolor}

% --- Styled callout boxes ---
\definecolor{keyborder}{RGB}{66,133,244}
\definecolor{keybg}{RGB}{240,248,255}

\newsavebox{\keyboxcontent}
\newenvironment{keybox}{%
  \par\smallskip\noindent
  \begin{lrbox}{\keyboxcontent}
  \begin{minipage}{\dimexpr\linewidth-2\fboxsep-2\fboxrule-1.5pt\relax}
  \raggedright
  \emergencystretch=1em
  {\color{keyborder}\faLightbulbO}\ \textbf{\textcolor{keyborder}{Observation.}}\space
}{%
  \end{minipage}
  \end{lrbox}
  \fcolorbox{keyborder}{keybg}{\usebox{\keyboxcontent}}%
  \par\smallskip
}

% --- End styled callout boxes ---

% --- FontAwesome icons for figures ---
\usepackage{fontawesome}


\hyphenpenalty=9000
\exhyphenpenalty=9000
\brokenpenalty=7000
\doublehyphendemerits=100000
\finalhyphendemerits=100000
\emergencystretch=4em


\begin{document}

\title{Beyond Storage: State as a Runtime Control Problem in Parallel and Distributed Systems}

\author{Shuhao Zhang}
\authornote{Corresponding author.}
\authornote{Supported by New Generation Artificial Intelligence-National Science and Technology Major Project (No. 2025ZD0123804).}
\orcid{0000-0002-9927-6925}
\affiliation{%
  \institution{Huazhong University of Science and Technology}
  \city{Wuhan}
  \country{China}}
\email{shuhaozhang@hust.edu.cn}

\author{Haoran Peng}
\orcid{0009-0003-6745-1762}
\affiliation{%
  \institution{Huazhong University of Science and Technology}
  \city{Wuhan}
  \country{China}}
\email{q1293132751@163.com}

\author{Chun Ho Ma}
\orcid{0009-0008-1636-0381}
\affiliation{%
  \institution{Huazhong University of Science and Technology}
  \city{Wuhan}
  \country{China}}
\email{u202490042@hust.edu.cn}

\author{Yancan Mao}
\orcid{0000-0002-7824-5978}
\affiliation{%
  \institution{ByteDance}
  \city{Singapore}
  \country{Singapore}}
\email{maoyancan@u.nus.edu}

\author{Yufeng Du}
\orcid{0009-0009-8919-0018}
\affiliation{%
  \institution{Huazhong University of Science and Technology}
  \city{Wuhan}
  \country{China}}
\email{yufeng050601@gmail.com}

\author{Shifeng Liu}
\orcid{0009-0001-1457-5839}
\affiliation{%
  \institution{Huazhong University of Science and Technology}
  \city{Wuhan}
  \country{China}}
\email{u202315663@hust.edu.cn}

\author{Ruijie Qiu}
\orcid{0009-0004-0659-3510}
\affiliation{%
  \institution{Huazhong University of Science and Technology}
  \city{Wuhan}
  \country{China}}
\email{Ying01020@outlook.com}

\author{Xiaofei Liao}
\orcid{0000-0002-2960-4982}
\affiliation{%
  \institution{Huazhong University of Science and Technology}
  \city{Wuhan}
  \country{China}}
\email{xfliao@hust.edu.cn}

\author{Hai Jin}
\orcid{0000-0002-3934-7605}
\affiliation{%
  \institution{Huazhong University of Science and Technology}
  \city{Wuhan}
  \country{China}}
\email{hjin@hust.edu.cn}

\begin{abstract}
Moving a hot shard, reclaiming a cached prefix, or repairing a vector index can improve one part of a service while increasing the cost of another. We survey state-management mechanisms in stream processing, analytics, large language model serving, retrieval, and continual learning. We trace each decision across three dimensions: access and scheduling, state-aware execution, and evolution and reuse. Three conditions explain when a mechanism can transfer across settings: which operations are valid for the state object, whether the action competes with foreground work, and whether reclaimed state can be recovered when it is needed. These comparisons motivate runtime contracts for authority, resource reservations, and later reconstruction. We apply them to saturation-aware migration, derive a conditional backlog bound, and propose tests that follow decisions through disturbance and recovery.
\end{abstract}

\ccsdesc[500]{Computer systems organization~Distributed architectures}
\ccsdesc[300]{Software and its engineering~Middleware}
\ccsdesc[300]{Information systems~Data management systems}
\ccsdesc[200]{Computing methodologies~Machine learning}
\keywords{state management, parallel and distributed systems, stream processing, LLM serving, retrieval-augmented generation, continual learning}
\maketitle
\renewcommand{\shortauthors}{Zhang et al.}

\section{Introduction}

A stream processor can relieve a busy worker by moving keys to another machine. The move saves work at the source after handoff, but the copy can further slow the events already queued there. An inference worker faces a similar decision when it reclaims a completed session's cached prefix to admit a waiting request. Admission becomes faster now, while the returning session may have to repeat prefill. In both settings, the first metric to improve can hide the cost shifted to another worker or a later request.

We ask when a local state decision produces a service-wide benefit. Streaming systems make progress and update authority explicit~\cite{akidau2013millwheel,murray2013naiad,zaharia2013dstreams,akidau2015dataflow}. Serving and retrieval systems add cases in which state may be replicated, dropped, or reconstructed~\cite{kwon2023pagedattention,zheng2023sglang,zhang2026flowrag}. A stream shard can move only if the system preserves the committed update boundary. A cached prefix can be rebuilt when the same inputs and compatible model state are available. The comparison explains which constraints a policy must carry when it transfers to a different state object.

We trace decisions in three dimensions. \emph{Access and scheduling} identifies who can use the object and where work waits. \emph{State-aware execution} follows the cost of layout, representation, and movement. \emph{Evolution and reuse} examines what survives for later work. Figure~\ref{fig:propagation} summarizes their interaction.

The survey contributes a common set of comparison questions, evidence about when mechanisms transfer, and design constraints that follow from those comparisons. Section~\ref{sec:foundations} defines the scope and questions. Section~\ref{sec:core-dimensions} traces how mechanisms work, and Section~\ref{sec:comparative-synthesis} compares their results under different conditions. Sections~\ref{sec:design-implications} and~\ref{sec:evaluation-outlook} derive runtime requirements and propose tests that measure what happens after an action.

\begin{figure*}[t]
\centering
\includegraphics{figures/pdf/state_propagation.pdf}
\caption{Feedback among access, execution, and reuse. The location of state changes the work needed to access it, and execution determines which state is available to later requests.}
\label{fig:propagation}
\end{figure*}

\section{Foundations: Scope, Model, and Taxonomy}
\label{sec:foundations}

We include studies that name a persistent state object, expose a runtime decision over it, and connect that decision to a service outcome such as latency, recovery, or result quality. The evidence covers streaming and transactional systems, hardware-conscious and approximate analytics, LLM serving, retrieval-augmented generation (RAG), and continual learning. Studies that explain static storage layouts or model updates provide context for these decisions.

\subsection{Survey Protocol and Evidence Posture}

We follow citations backward to the origin of a mechanism and forward to later implementations and evaluations. We prefer published versions where available and select cases that expose differing state, decisions, or outcomes. This is a representative comparison. For each study we keep the workload, baseline, and measurement window with the claim. For example, a shorter migration pause measures the handoff, while foreground recovery requires measurements after the change in demand.

\subsection{Positioning Relative to Existing Surveys}

Prior surveys describe stream processing, transactional and hardware-conscious execution, and continual learning~\cite{hirzel2014catalog,cugola2012processing,zhang2024survey,zhang2020hardware,delange2021survey}. Our unit of comparison is a runtime decision over state. This lets us compare, for example, when a stream worker can move a shard with when an inference worker can reclaim a reusable prefix. Storage and cache surveys supply adjacent mechanisms such as compaction and eviction~\cite{luo2020lsm,ali2011survey}. We use them where those decisions affect later execution or recovery.

\subsection{Broader Literature Landscape}

Appendix~A of the supplement traces the literature behind the comparison. It covers progress and authority in dataflow, placement and transfer in serving and tiered memory, and the effect of maintenance and retention on later service.

\subsection{System Model: Where State Lives in a Runtime}

An object qualifies as managed state when it survives one operation, affects later work, and is subject to a runtime decision. A page in the key--value (KV) cache qualifies because a later decoding step can read it and the allocator can place, move, or reclaim it. Table~\ref{tab:state-system-model} shows where state lives and how it moves in representative runtimes.

\begin{table}[!htbp]
\centering
\caption{State in representative runtime architectures.}
\label{tab:state-system-model}
\small
\setlength{\tabcolsep}{4pt}
\renewcommand{\arraystretch}{1.10}
\resizebox{\textwidth}{!}{%
\begin{tabular}{@{}p{3.1cm}p{3.1cm}p{4.2cm}p{6.2cm}@{}}
\toprule
\textbf{System class} &
\textbf{Computational unit} &
\textbf{Communication path} &
\textbf{Representative state objects} \\
\midrule
Stream processing &
Operators and tasks &
Streams, shuffles, backpressure channels &
Keyed state, windows, operator state, watermarks, checkpoint metadata \\
% \addlinespace
Transactional dataflow and databases &
Transactions, execution fragments, workers &
Dependency exchange, log replication, remote reads &
Records, indexes, versions, locks, dependency graphs, logs, recovery snapshots \\
% \addlinespace
Distributed task/actor runtimes &
Tasks, actors, executors &
RPCs, object-store transfers, lineage replay &
Objects, actor-local state, lineage metadata, placement state, scheduler queues \\
% \addlinespace
LLM serving &
Prefill and decode workers &
KV transfer, model-parallel communication, request queues &
KV caches, page tables, prefix trees, adapter state, model-residency metadata \\
% \addlinespace
RAG and agentic workflows &
Retrievers, planners, tools, generators &
Vector-index queries, tool calls, memory updates &
Vector indexes, document memories, workflow context, tool traces, retained memories \\
% \addlinespace
Edge and approximate analytics &
CPU/GPU operators, sensor tasks &
Device transfer, edge--cloud transfer &
Summaries, sketches, compression dictionaries, quality-control metadata \\
\bottomrule
\end{tabular}%
}
\end{table}


An object may have a logical identity, a physical representation, and metadata that controls its use. For a stream shard, key ranges identify the state, a snapshot encodes its contents, and the ownership record and input position determine where execution resumes. A move changes placement and authority together. Keeping these roles distinct helps identify which records must change atomically and which can be reconstructed.

\subsection{A Unifying View of Stateful Computing Systems}

\subsubsection{Who Shares the State?}

Two requests can read the same cached prefix. If a worker reclaims its last copy to admit one request, the next request may have to rebuild it. For a stream shard, many records can update the same key. The runtime must establish which worker may apply each update.
\subsubsection{The Propagation Perspective}

Consider a request that needs state held by another worker. Its placement determines how far the data must travel. Copying that data uses bandwidth that other requests may already need. If the destination keeps a valid copy, a later request may avoid the trip. These are the access, execution, and reuse effects depicted in Figure~\ref{fig:propagation}. The same dependencies arise in operator placement, approximate execution, and online retention~\cite{zhang2019briskstream,zeng2024pecj,zhou2025ferret}.

Suppose a worker serves requests in arrival order, one at a time. A new request arrives behind $q$ requests, each with $s$ milliseconds of service remaining. If its own first-token work takes $t$, then its time to first token (TTFT) is $qs+t$. Two preceding requests at 5~ms each add 10~ms. The finite response $\Delta\mathrm{TTFT}/\Delta q$ is then 5~ms per additional queued request. At fixed $q$, a cache miss adds 4~ms if recomputing the prefix takes 4~ms. For an empirical estimate, use paired replays and hold the other factors fixed. Continuous batching changes service overlap, so the queue response depends on the operating point.

\begin{keybox}
Compare an action over the requests that benefit immediately and the later requests that absorb its cost.
\end{keybox}

\subsection{Taxonomy of State Management Problems}

Table~\ref{tab:taxonomy} groups decisions by the problem they address. One system can exhibit all three dimensions: a serving engine can route around an overloaded worker, allocate memory for active requests, and decide which prefixes to keep for later reuse.

\begin{table*}[t]
\small
\caption{Taxonomy of state-management problems.}
\label{tab:taxonomy}
\begin{tabular}{@{}p{0.18\textwidth}p{0.37\textwidth}p{0.37\textwidth}@{}}
\toprule
Dimension & Core question & Representative concerns \\
\midrule
Access and scheduling & How is shared state exposed, costed, and scheduled under concurrency? & hotspot diagnosis, conflict propagation, locality, topology, recovery, migration, prefix-aware sharing \\

Execution optimization & How does state interact with hardware and service objectives? & placement, compression, KV-cache paging, prefill/decode coupling, approximation, latency-energy-quality trade-offs \\

Evolution and reuse & How is state updated, retained, and reused over time? & online updates, memory budgets, cache growth and reclamation, sample selection, retriever drift, knowledge maintenance \\
\bottomrule
\end{tabular}
\end{table*}

\subsubsection{Five Questions for Comparing State Decisions}

Access, execution, and evolution organize Section~\ref{sec:core-dimensions}. The five questions in Table~\ref{tab:scaffold} are used within each topic. Consider a cached prefix (the object) that the allocator can reclaim (the decision). Reclamation admits a waiting request but may force a returning user to repeat prefill (the effect on later work). Time to first token across both visits captures the tradeoff (the measurement). If the study ends before the user returns, how much delay reclamation caused on that visit remains open.

\begin{table*}[t]
\small
\caption{Questions used to compare state-management mechanisms.}
\label{tab:scaffold}
\begin{tabular}{@{}p{0.18\textwidth}p{0.37\textwidth}p{0.37\textwidth}@{}}
\toprule
Question & What to extract from each paper & Typical answers in this survey \\
\midrule
State object & Which runtime object persists across tasks, requests, or time? & windows, indexes, logs, dictionaries, KV caches, retriever memories, coresets \\

Control surface & What decision does the system make over that object? & placement, batching, migration, paging, admission, retention, reuse orchestration \\

Coupling path & Why does this state decision propagate to end-to-end behavior? & contention, locality, memory pressure, phase asymmetry, drift, quality-loss amplification \\

Evaluation boundary & Which system-level property is optimized or bounded? & throughput, p99 latency, recovery time, energy, bounded error, forgetting, reuse quality \\

Unresolved contract & What is still missing after the paper's contribution? & weak observability, no closed loop, limited portability, incomplete lifecycle control, no long-horizon evaluation \\
\bottomrule
\end{tabular}
\end{table*}

\subsubsection{Representative System Classes}

The comparisons use five settings. Dataflow foregrounds authority and recovery, analytics foregrounds hardware cost and approximation error, and serving foregrounds queueing and memory pressure. Retrieval adds update visibility, while continual learning tests the longer-term effect of losing retained data.

\section{Core Dimensions of State Management}
\label{sec:core-dimensions}

Three examples recur in this section: (A) a streaming fraud detector with skewed account traffic, (B) a multi-tenant LLM assistant under KV-memory pressure, and (C) a RAG service whose corpus changes while queries arrive. Each starts with one state decision and follows the work it creates elsewhere.

\subsection{State Access and Scheduling}

\subsubsection{From Invisible Contention to Observable State Access}

In running example A, a fraud detector updates the recent transaction history of each account. Suppose several independently updated account keys are assigned to the same worker. After a marketing campaign starts, events for these accounts arrive faster than this worker can process them. A second worker has spare capacity and could take over some of the accounts. A single busy account would still require ordered updates after a move. Per-key arrival rates and the queue at each worker distinguish that case from several independent accounts competing for one worker. Queueing and synchronization can limit stream throughput even when the operator code itself scales~\cite{zhang2017revisiting,zhang2017mqo}.

Early adaptive stream systems made placement part of the runtime's decision~\cite{abadi2003aurora,abadi2005borealis}. Progress tracking in MillWheel, Naiad, and Dataflow provides a different prerequisite: it tells a worker when earlier work has been accounted for~\cite{akidau2013millwheel,murray2013naiad,akidau2015dataflow}. A reactive controller needs both an indication of where the load is concentrated and a progress boundary at which it can move the state.

\subsubsection{Cost Modeling Under Locality and Topology}

In Example~A, the candidate move assigns some of the busy accounts to a second worker. Estimate waiting time under both placements for the same incoming events, including those already served by the destination. After the move, each worker has a different arrival rate. Its processing rate may change as well: BriskStream models how placing communicating operators on different sockets changes execution time on a shared-memory machine~\cite{zhang2019briskstream}. Between workers there is also a one-time cost before the new placement can take effect.

While writes to the moving keys are fenced, their events must wait. Snapshotting and serialization consume CPU and memory bandwidth at the source. The copy reads from source memory, crosses the network, and writes to destination memory. The available rate of the slowest resource can limit the copy even if the destination has idle CPU. Megaphone shows why granularity matters: batching a planned migration reduced the disruption for several stateful NEXMark queries, with results varying by the amount of state moved~\cite{nasir2019megaphone}.

Before moving write authority, consider redirecting new work, reading an existing valid replica, or deferring the action if those options preserve the service contract. Compare total waiting through the burst and queue drain, including transfer work at the source. A short burst may end before shorter queues can repay the handoff cost. Even when a move is expected to pay off, the source, destination and link need spare capacity for the first batch. If a required resource is saturated, defer the move before fencing writes or taking a snapshot. Section~\ref{sec:migration-admission} gives the checks for starting, pausing and resuming.

\subsubsection{Runtime Control, Recovery, and Stateful Governance}

Before fencing, reserve headroom at the source, destination and shared link, and check the migration budget and recovery epoch (Section~\ref{sec:migration-admission}). Once those conditions hold, the next question is where processing resumes. The old and new owners must agree on the last event processed for each key before the new owner accepts updates. Checkpoint and replay protocols record the input position from which processing can resume~\cite{carbone2015asynchronous,zhao2024recovery}.

A handoff has four responsibilities: drain or fence writes at a known boundary, transfer the shard and its progress record, commit new authority, and resume from that record. Megaphone and Spacker expose different parts of the move, while checkpointing and fast-recovery systems show how to retain the position needed for replay~\cite{nasir2019megaphone,mao2025spacker,carbone2015asynchronous,zhao2024recovery}. A committed ownership record tells recovery which worker is authoritative. Replay and sink semantics determine how duplicate processing affects visible output.

Other controllers can use the same resource during the handoff. For example, a checkpoint may take a snapshot while a migration serializes the same shard. Coordination must account for that concurrency before either controller reserves capacity.

\subsubsection{Open Challenges in Access Management}

The access problem has two independent checks. A runtime needs enough information to tell a hot key from a hot placement, and it needs a way to preserve progress while relieving the contention. Exposing finer-grained contention helps choose the object to move. It leaves the resource and ownership questions open whenever migration, recovery, and foreground admission run under separate controllers.

\subsection{State-Aware Execution Optimization}

\subsubsection{Placement and Movement}

The stateful data path runs from the worker that holds the data to the operation that uses it. FineStream can divide window work between CPU and GPU under a shared-memory topology~\cite{zhang2020finestream}. For a discrete accelerator, moving the same window can add a link traversal; for co-running operators, even shared memory can become the contested resource~\cite{zhang2016corun,zhang2016elastic}. A local speedup matters when it exceeds routing, transfer, and contention costs on the affected path. The measurement must follow the result through the same service boundary under both placements.

\subsubsection{Granularity, Scheduling, and Representation}

Example~B exposes three different causes of an inference request's wait. PagedAttention changes the number of active sequences that can fit by allocating KV blocks as sequences grow~\cite{kwon2023pagedattention}. Sarathi-Serve changes when a decoding sequence can run by splitting prefills into smaller units~\cite{agrawal2024sarathi}. DistServe isolates the two phases on separate workers and adds a KV transfer between them~\cite{zhong2024distserve}. Measuring first-token latency and inter-token latency separately reveals whether an action removed a wait or moved it to another phase. A useful comparison keeps the prompt and output-length distributions, and the latency targets, explicit.

Representation changes which costs can be exchanged. Offloading exchanges GPU capacity for host memory and link traffic. Lossless compression exchanges bytes for encoding and decoding work, as illustrated by the stateful dictionaries in CStream~\cite{zeng2023cstreamicde,zeng2024cstream}. Dropping context or approximating a join can also alter the returned result. PECJ makes late and missing records part of the output error model, while LibAMM bounds a different quantity, numerical approximation error~\cite{zeng2024pecj,zeng2024libamm}. The open scheduling question is how to choose among these actions under the current bottleneck and the applicable output-quality constraint.

\subsection{State Evolution and Reuse}

\subsubsection{Update Visibility and Maintenance}

In Example~C, a knowledge service accepts new or deleted documents while serving retrieval queries. It must decide when an update becomes visible and when to spend resources on maintenance. Searching a temporary index exposes an insert before a merge, as FreshDiskANN illustrates~\cite{singh2021freshdiskann}. This shifts some lookup and repair work into later requests. A retrieval trigger such as Self-RAG chooses when to consult memory and operates on whatever version was published~\cite{asai2024selfrag}. The two decisions are coupled through which evidence a query can reach. Measuring recall on a fixed corpus hides both publication lag and the scheduling cost of a concurrent repair~\cite{wang2026candor}.

\subsubsection{Retention and Reconstruction}

When the service reclaims state, later work can use a surviving copy, rebuild the object, or continue without it if semantics permit. Serving can rebuild a compatible prefix from its tokens and the matching model configuration. Index repair can recover a representation if the document remains available; deleting the document itself has a different visibility requirement. Continual learning faces another case: if no upstream copy exists, removal of the only retained exemplar makes exact replay of that example impossible. Policies that adapt sample selection to memory budgets must then be compared on later task accuracy and class coverage as the stream evolves~\cite{zhou2025ferret,li2026streamfp}.

The control decision depends on the time until the next use, the cost of restoring state, and the harm if it cannot be recovered. Observe misses and recovery work following an eviction, query behavior across publication, and task quality after subsequent training. Each measurement window must cover the later work that the action can change.

\subsection{How KV Eviction Affects the Next Request}
\label{sec:propagation-trace}

Session A has finished its first request, and the worker retains its unpinned KV prefix. Request B arrives with a long prompt and can start only after enough pages become free. The allocator can evict A's prefix or wait. If A returns while its prefix would still have been cached, eviction can force a reload or recomputation~\cite{kwon2023pagedattention,hu2024memserve,liu2025lmcache}.

Assume retaining A's prefix adds 10~ms to B's first-token latency. Eviction removes that wait and has negligible direct cost. If A returns with probability \(p\) before it would have been evicted anyway, and there is at most one return, eviction changes summed TTFT by \(-10+p d\)~ms per episode. Here \(d\) is the additional latency on A's visit. A 4~ms recomputation penalty favors eviction for every \(p\). A 20~ms penalty favors retention when \(p>0.5\). Both policies must satisfy the per-request latency budget. All inputs in this example are hypothetical.

A valid remote copy can change the result. MemServe compares transfer time with the avoided prefill cost~\cite[Sec.~5.3.1]{hu2024memserve}. Using that decision rule, suppose a remote copy needs an assumed 16~MB transfer and 2~ms of setup, with transfer and setup both on the critical path. At an effective 4~GB/s (decimal units), the reload takes 6~ms. At 0.8~GB/s it takes 22~ms, so a 20~ms recomputation is preferable. The crossover occurs at about 0.89~GB/s. With \(p=0.75\), the faster path changes the expected cost of eviction from a 5~ms increase at the slower link rate to a 5.5~ms decrease at the faster one. Only the decision rule comes from MemServe.

The example traces one operation through all three dimensions: admission improves for B, transfer or prefill increases for A, and the retained copy changes the result. Section~\ref{sec:comparative-synthesis} asks which parts of this reasoning survive when the state governs writes or cannot be rebuilt.

\section{Comparative Synthesis Across Domains}
\label{sec:comparative-synthesis}

Section~\ref{sec:core-dimensions} shows how a decision propagates. A comparison must also explain why the same move is appropriate in one system and unsuitable in another. We ask three questions. What makes use of the state valid? Which foreground work competes with the action? Can that state be recovered before reuse? These conditions distinguish the evidence a system already provides from the test needed to transfer its result.

\subsection{When Can State Be Used, and When Is It Useful?}

Validity is a property of an operation on a particular state. Spanner orders committed transactions consistently with real time, and a stream worker must hold write authority before applying an update~\cite{corbett2012spanner,nasir2019megaphone}. A retrieval service can check whether a replica exposes the requested version. The usefulness of that valid state depends on the workload: an index can serve the correct version and still miss relevant documents. Measure usefulness through recall over a stated query distribution or earlier-task accuracy after updates. Those measurements complement the version and authority checks.

\subsection{Cross-Domain Comparative Synthesis}

Use Table~\ref{tab:comparison-guide} to identify the condition that changes the decision in each domain. Table~\ref{tab:representative-mechanism-matrix} then pairs each mechanism with evidence and a test for extending the result. This separates what a system has demonstrated from what an analogous deployment would still need to establish.

\begin{table*}[t]
\small
\centering
\setlength{\tabcolsep}{4pt}
\caption{Conditions that change a state-management decision.}
\label{tab:comparison-guide}
\begin{tabular}{@{}p{0.17\textwidth}p{0.37\textwidth}p{0.39\textwidth}@{}}
\toprule
Domain & Condition that changes the decision & Action supported by the cited mechanism \\
\midrule
Streaming & Is a worker saving a recovery point or handing a key to a new owner? & Restore from a checkpoint, or move keys at a progress boundary~\cite{carbone2015flink,nasir2019megaphone}. \\
Analytics & Must the query use all records, or will the user accept an estimate? & Partition exact window processing, or select a sample under a time or error constraint~\cite{zhang2020finestream,agarwal2013blinkdb}. \\
LLM serving & How much of a valid prefix is local, and how long would fetching the rest take? & Route to a worker, then fetch missing blocks or recompute them~\cite{qin2025mooncake,hu2024memserve}. \\
Vector search & Are recent inserts kept in a temporary graph or written to a partition of the main index? & Merge the temporary index, or split and repair an affected partition~\cite{singh2021freshdiskann,xu2023spfresh}. \\
Continual learning & Must the model adapt as data arrive, or can training wait until evaluation? & Select replay examples for this update, or balance the buffer for later training~\cite{aljundi2019mir,prabhu2020gdumb}. \\
\bottomrule
\end{tabular}
\end{table*}

\subsubsection{Streaming and Transactional Dataflow Systems}

Flink and Megaphone both need a boundary between updates already represented in state and updates still to be applied. In Flink's aligned checkpointing protocol, barriers separate the records whose effects belong in the snapshot from later input~\cite{carbone2015flink,carbone2015asynchronous}. Recovery restores the checkpoint and replays later input. Megaphone instead moves state during live execution. A timestamp frontier identifies when the old worker has accounted for earlier records and the new worker can take over the corresponding keys~\cite{nasir2019megaphone}. A live reconfiguration therefore needs both a progress boundary and capacity to perform the handoff at that boundary.

\subsubsection{Hardware-Conscious and Approximation-Aware Stateful Execution}

FineStream~\cite{zhang2020finestream} runs window queries on a chip where the CPU and GPU share memory. To finish a window sooner, it divides the work between the two processors while keeping the same records in the computation. BlinkDB~\cite{agarwal2013blinkdb} serves analytical queries whose users can accept an estimate. Given a time limit or an error bound, it chooses how large a sample to query and reports the estimate with its uncertainty. Moving work to the GPU changes when a window result is ready. A smaller sample can also change the reported value. A latency result transfers to approximate execution only when evaluated against the output-quality requirement.

\subsubsection{LLM Serving and Short-Lived Memory Lifecycles}\label{sec:serving-comparison}

Suppose the worker with the longest matching prefix also has the longest queue. Mooncake~\cite{qin2025mooncake} can route to another worker if its queue, local cache and fetch cost imply an earlier first token. MemServe~\cite{hu2024memserve} also considers queues and cache hits. After routing, its cost-based policy fetches additional blocks when the estimated transfer takes no longer than the prefill it would save. The estimate uses the maximum transfer bandwidth, and its evaluation used eight GPUs in one server. Under inter-server contention, the decision depends on the bandwidth available during the copy.

LMCache exposes cache lookup, pinning, and movement through its controller, and connects these operations to inference engines through a connector~\cite{liu2025lmcache}. Together with MemServe's distributed pool and scheduling policy, these interfaces let the runtime select between a surviving copy and recomputation. Safe reuse still requires matching the model, adapter, and prefix context. A completed prefix can have multiple readers. In contrast, Llumnix hands off a running request and must carry KV entries created during the copy so that decoding can continue~\cite{sun2024llumnix}. Neither a lookup result nor a move primitive reserves the shared link. Evaluations must account for retrieval and foreground latency when multiple transfers compete.

\subsubsection{Retrieval Memory and Vector-Index State Layers}

FreshDiskANN~\cite{singh2021freshdiskann} keeps its main search graph on SSD. It puts new vectors in an in-memory index and searches both indexes, filtering deleted entries from the results. Queries can find an insertion before the next merge. When the temporary indexes reach a memory threshold, a background merge incorporates the updates into the disk graph. SPFresh~\cite{xu2023spfresh} partitions the vectors by proximity. It inserts a new vector into the nearest partition. If the partition grows too large, it splits it and reassigns affected vectors near the new boundary. Thus a burst of insertions leaves different work behind: a merge with the disk graph in FreshDiskANN, or repairs to the affected partitions in SPFresh. To transfer a result across the two policies, align the visibility of new versions and measure recall and latency during repair.

\subsubsection{Continual Learning and Retention Governance}

MIR~\cite{aljundi2019mir} updates the model as new batches arrive. It estimates how a step on the new batch would change the loss on stored examples. It replays the examples with the largest predicted increases alongside the new batch. GDumb~\cite{prabhu2020gdumb} can defer training until evaluation. As examples arrive, it keeps a buffer whose class counts it tries to balance. At evaluation it trains a fresh model using only that buffer. For MIR, the immediate question is which old example to replay in this update. In GDumb, losing the last example of a class can remove it from the next training run. The benefit of retention must therefore be measured on future tasks as well as the current update.

\subsubsection{Comparing the Effects of State Decisions}

Progress boundaries let recovery and migration identify the state to restore or move. They also determine whose work must wait. Splitting a transition only helps where the dependence itself can be split. Megaphone can move independent groups while other keys progress. Sarathi-Serve can schedule decode work between prefill chunks. Neither technique transfers to a single ordered hot key or to an execution phase that cannot be interrupted. Table~\ref{tab:representative-mechanism-matrix} relates each mechanism to the result observed and the remaining transfer test.

\begin{table*}[t]
\small
\centering
\setlength{\tabcolsep}{3pt}
\caption{State decisions in the surveyed systems and their consequences for subsequent work.}
\label{tab:representative-mechanism-matrix}
\begin{tabular}{@{}p{0.15\textwidth}p{0.40\textwidth}p{0.39\textwidth}@{}}
\toprule
Decision & Evidence from the cited work & Open test and its reason \\
\midrule
Stream progress & MillWheel records keyed state and processing progress so computation can resume after failure~\cite{akidau2013millwheel}. & Interrupt a move on either side of the commit to check authority, replay, and duplicate effects. \\
Shard migration & In Megaphone's planned NEXMark Q4 move, the latency spike exceeded 2~s all at once, versus an increase of at most 100~ms with batching~\cite[Figure~8]{nasir2019megaphone}. & Measure foreground backlog when skew triggers the move and traffic competes for the copy resources. \\
Device placement & FineStream divides window processing between a CPU and GPU sharing memory~\cite{zhang2020finestream}. & On discrete devices, include the cost of moving intermediate results across the link. \\
Approximate queries & BlinkDB selects a sample under a response-time or error constraint and reports the uncertainty~\cite{agarwal2013blinkdb}. & If the estimate triggers an alert, its uncertainty matters most near the alert threshold. \\
Serving capacity & PagedAttention allocates KV blocks on demand. Sarathi-Serve interleaves chunks of prefill with decoding~\cite{kwon2023pagedattention,agrawal2024sarathi}. & Measure admission and inter-token latency together when long prompts arrive during decoding. \\
Prefix transfer & Mooncake accounts for queues and cache hits. MemServe compares transfer with recomputation. LMCache supplies block lookup and movement~\cite{qin2025mooncake,hu2024memserve,liu2025lmcache}. & Measure return-visit TTFT under shared-link pressure because even a cache hit can incur transfer and queue delay. \\
Adapter residency & S-LoRA manages adapter weights and KV tensors in one memory pool~\cite{sheng2023slora}. & After switching adapters, check that existing KV blocks still match the parameters used for this request. \\
Index maintenance & At one billion vectors and 1\% daily updates, SPFresh reports lower peak maintenance resources than a global-rebuild baseline~\cite{xu2023spfresh}. FreshDiskANN searches temporary and disk indexes~\cite{singh2021freshdiskann}. & Measure recall, visibility lag, and query latency when changes cluster near partition boundaries. \\
On-demand retrieval & Self-RAG decides during generation whether to look up documents~\cite{asai2024selfrag}. & Hold the lookup policy fixed and vary index visibility to separate missed updates from missed retrievals. \\
Retained examples & MIR replays examples at risk in the next update. GDumb trains at evaluation on a class-balanced buffer~\cite{aljundi2019mir,prabhu2020gdumb}. & After shrinking the buffer, test accuracy by earlier task and class to locate effects of losing coverage. \\
\bottomrule
\end{tabular}
\end{table*}

\subsection{What the Measurements Establish}

\paragraph{Validity follows the object.}
A progress frontier identifies the updates a new owner needs to apply. The ownership and replay rules determine whether those updates take effect once. A compatible prefix can have many readers, and a valid retrieval version may still answer poorly on a shifted corpus. A shared abstraction needs to carry the appropriate authority or version check together with the distinct measure of expected usefulness.

\paragraph{Resource interference follows the path.}
The Megaphone comparison isolates the grain of a scheduled reconfiguration. It shows that smaller moves can lower the intervention peak under those conditions. A reactive move also depends on the resources available at the instant it starts. MemServe uses the same reasoning about local work versus transfer cost for a read path. Its single-server result and its transfer cost model leave cross-server contention as a separate test. Both cases motivate reserving capacity where the action runs.

\paragraph{Reconstruction determines the evaluation horizon.}
A compatible cached prefix can be rebuilt at a later latency cost. An evicted exemplar without another source cannot be replayed. When the cost appears on later requests or tasks, end the observation after revisit or after the relevant updates. These three consequences lead to the design requirements in Section~\ref{sec:design-implications}: preserve validity, reserve shared resources, and account for later effects.

\section{Design Implications for Stateful Runtimes}
\label{sec:design-implications}
The comparisons yield three requirements. The runtime must record who may use or change the state, reserve the resources consumed by the action, and identify who pays the cost if state must later be recovered. A destination may be eligible to own a shard even when the source lacks the bandwidth to copy it. The architecture has to enforce both checks at the point of handoff.

\subsection{A Design Space for Stateful Runtime Architectures}

\subsubsection{State Granularity: Record, Key, Segment, and Session}

The object that can safely move sets the smallest useful unit of control. Key groups with independent updates can move separately. KV pages allow fine-grained allocation, while whole prefixes and sessions can simplify reuse and accounting~\cite{zhang2019briskstream,kwon2023pagedattention,zheng2023sglang}. Smaller units give the scheduler more options but require more metadata and coordination. Larger units simplify control at the cost of moving or retaining state that the next operation may never use.

\subsubsection{Control-Plane Coupling: Embedded, Sidecar, or Shared Service}

An embedded scheduler can act on a local queue immediately. A sidecar can coordinate several workers, but the queue may change between observation and actuation. A shared service can serialize reservations only if its interface orders competing requests against the same resource budget. In all three designs, whoever starts a transfer must check that the reservation is still valid (Section~\ref{sec:migration-admission}).

\subsubsection{Lifecycle Management and Visibility}

A runtime admits, places, updates, shares or transfers, compacts, and eventually reclaims state. A stream processor transfers a shard and resumes processing at its new owner. A serving engine admits a request and later reclaims its KV pages. A retrieval service updates or compacts an index as its corpus changes. For each action, the runtime needs to record which requests can still use the object afterwards.

\subsubsection{Service-Boundary Alignment: Local Wins vs End-to-End Wins}

Set the objective at the point where the user observes service. Faster prefill can delay active decoding. Faster publication can increase the repair work that queries share. Reducing a replay buffer can harm later task accuracy. A controller needs a metric for its immediate benefit and a budget or invariant for the cost borne elsewhere.

\subsection{Failure Modes and Anti-Patterns}

Table~\ref{tab:antipattern-repairs} pairs observable failures with the control rule needed to prevent or detect them. Measure foreground completions as well as the resource targeted by a repair.

\begin{table}[t]
\small
\caption{Recurring anti-patterns and the contract repairs that make them reviewable.}
\label{tab:antipattern-repairs}
\begin{tabular}{@{}p{0.22\textwidth}p{0.26\textwidth}p{0.18\textwidth}p{0.26\textwidth}@{}}
\toprule
Anti-pattern & Missing contract & Primary blueprint layer & Expected design effect \\
\midrule
Stale or ambiguous telemetry & Identify the object, sample age, and admissible action & Observation & Defer an action when the observation no longer supports it. \\

Conflicting policies & Define which action takes precedence & Decision and admission & Preserve foreground capacity when actions compete. \\

Steady-state-only evaluation & Define the disturbance and recovery windows & Observation and recovery & Make delayed effects visible after the perturbation. \\

Uncharged maintenance & Account for update, compaction, and deletion work alongside read gains & Catalog and actuation & Record the work queued by a faster foreground path. \\

Copying at saturation & Charge copy work against the foreground bottleneck & Decision and admission & Defer before fencing if resource headroom is unavailable. \\

Concurrent migration to a shared bottleneck & Reserve shared links and destination capacity together & Actuation and coordination & Serialize conflicting moves and back off before retrying. \\

Stale write authority & Preserve ownership and progress across migration and recovery & Actuation and recovery & Reopen writes only at the committed owner. \\
\bottomrule
\end{tabular}
\end{table}

\subsection{A Contract-Oriented Blueprint for Stateful Runtimes}

Five interfaces make the requirements enforceable. A \emph{state catalog} names each object, its granularity, and who may read or change it. \emph{Observation} reports load, progress, and when each measurement was taken. \emph{Admission} decides whether a proposed action can meet the service limits and reserves shared capacity. \emph{Actuation} records which controller may execute the action and under which ownership epoch. \emph{Recovery} defines how to cancel before commit or resume after it. The following subsection shows these interfaces at work in the same migration decision.

\subsection{Saturation-Aware Migration}\label{sec:migration-admission}

\paragraph{How a move can amplify a burst.}
Fencing, snapshotting, copying, and replay can compete with foreground requests for CPU, memory bandwidth, storage I/O, and shared links. Near capacity, a small reduction in effective service rate can cause a much larger increase in waiting time. If a controller treats this queue growth as fresh evidence of skew, it may start more copies. Several independent controllers can then overload the same destination or network path. This is a possible source of the sustaining feedback described in studies of metastable failure~\cite{bronson2021metastable}. Megaphone illustrates how staging changes migration granularity~\cite{nasir2019megaphone}. Live VM migration motivates bounded pre-copy and explicit accounting for update churn~\cite{clark2005livemigration}. The admission check below addresses the foreground capacity needed for either kind of move.

\paragraph{Estimate the whole move.}
Let \(T_{\mathrm{mig}}\) be the time from starting a move to committing the handoff. When the phases are serial, \(T_{\mathrm{mig}}=T_{\mathrm{setup}}+B/b_{\mathrm{eff}}+T_{\mathrm{handoff}}\). Here \(B\) includes bulk and catch-up bytes, \(b_{\mathrm{eff}}\) is the transfer rate under foreground load, and setup and handoff times cover work outside the transfer interval. For overlapping phases, use the critical path. Bound the write pause separately, including any transfer performed inside the fence.

Compare movement with a reference policy on the same offered workload. Let \(C(x)\) be the extra cumulative foreground waiting from action \(x\) before settlement, in request-seconds, including waiting imposed by fencing and expected recovery work. Record rejected work separately. Index affected foreground classes and endpoints jointly by \(i\). With rates \(\lambda_i\) and mean waits \(W_i\) for the reference and primed quantities for the moved placement, the post-settlement saving per second is
\[
  D(x)=\sum_i \bigl(\lambda_i W_i-\lambda'_i W'_i\bigr).
\]

For an estimated remaining skew horizon \(H\) after detection, let \(T_{\mathrm{settle}}\) be the time to drain the migration transient. Require \(H>T_{\mathrm{mig}}+T_{\mathrm{settle}}\) and a conservative estimate of
\[
  G(x)=(H-T_{\mathrm{mig}}-T_{\mathrm{settle}})D(x)-C(x)>0.
\]
The horizon test excludes short bursts and the benefit test includes the wait added at the destination. Reject a move whose projected post-handoff service would violate either endpoint's limits. Compare waiting costs at a matched admission policy. When policies admit different work, also compare rejection and completion rates under the same offered load. Alternatives include deferral, redirecting new work, or reading a valid replica where semantics permit.

\paragraph{Reserve capacity and limit feedback.}
Let \(L_r\) bound foreground demand on resource \(r\), \(M_r\) bound admitted migration work, \(m_r\) bound the next batch, and \(H_r\) account for uncertainty. Express these as rates over the same control interval. Admit only when \(L_r+M_r+m_r+H_r\le C_r\) for every required resource, where \(C_r\) is a lower bound on usable capacity. Include both endpoints, the network path, and completion or abort work. A token bucket limits the transient copy burst; the shared reservation prevents multiple controllers from each assuming the same spare capacity. Check the write pause separately before fencing.

The rate check protects a share of resources. A bound on foreground backlog needs a bound on burst size and on the actual service delivered. Suppose admitted foreground arrivals in every interval of length \(t\) are at most \(\sigma+\rho t\) records, counting any backlog at the start of observation in \(\sigma\). Suppose that every continuously backlogged interval receives at least \(R(t-d)^+\) records of service, with \(z^+=\max(0,z)\), \(R>\rho\), and \(d\) covering the effect of repeated fences. This is a rate--latency service bound~\cite{leboudec2001networkcalculus}. Then the backlog is at most
\[
\sup_{t\ge 0}\{\sigma+\rho t-R(t-d)^+\}=\sigma+\rho d .
\]
Indeed, use the interval since the queue last emptied, or since observation began. The expression grows through \(t=d\) and decreases thereafter. Accounting for CPU or link headroom alone does not establish the service premise, especially for the fenced keys. A controller must supply a defensible \(R,d\) for each protected class and downstream bottleneck before treating the bound as a guarantee.

If observations are stale or the demand envelope is exceeded, suspend new batches before fencing and reduce interruptible copy work. Preserve resources for finishing or aborting a committed move. During a sustained overload, apply admission control or route work to a valid replica if the service contract permits it. Resume at a low rate only when measured foreground backlogs are draining and backpressure stays below a lower restart threshold for a recovery interval. This hysteresis prevents the admission rule from restarting a copy on a brief dip in pressure.

\begin{table}[t]
\small
\setlength{\tabcolsep}{4pt}
\caption{Contract fields checked when a runtime considers moving mutable state.}
\label{tab:migration-contract}
\begin{tabular}{@{}p{0.25\textwidth}p{0.38\textwidth}p{0.31\textwidth}@{}}
\toprule
Field & What it protects & Response when the check fails \\
\midrule
Saturation guard & Foreground capacity at both workers and along the shared path & Defer before taking a snapshot or fencing. \\
Migration budget & Per-resource rates and the cluster-wide in-flight limit & Reject or queue the reservation. \\
Negative-feedback boundary & Foreground queue growth and backpressure during a copy & Reduce interruptible copy rate, stop new batches, and apply restart hysteresis. \\
Rollback scope & Single write authority and the progress boundary for replay & Release a pre-commit fence at the old owner, or recover at the committed epoch. \\
Coordination scope & Shared reservations and exclusive destination leases for conflicting moves & Serialize conflicts and retry after randomized backoff. \\
\bottomrule
\end{tabular}
\end{table}

\paragraph{Coordinate across workers.}
Use shared per-resource token buckets and destination leases acquired before fencing. A lease carries an epoch so a delayed controller cannot authorize a second owner. Record completed directions for each object to detect reversals. A possible additional rule is to require a drop of at least \(\epsilon>0\) in a cluster score \(\Phi=\sum_i[(u_i-u_*)^2+\gamma(q_i/q_*)^2]\) after settlement. Here \(u_i\) is normalized bottleneck utilization, \(u_*\) a target, \(q_i\) foreground backlog, \(q_*>0\) a reference backlog, and \(\gamma\geq0\) a workload-specific weight. Compare settled windows under the same demand and account for all shared bottlenecks in the resource guard.

\emph{Conditional bound.} Consider successive settled configurations under fixed external demand and capacity. If each accepted move or coordinated group reduces the same lower-bounded score by at least \(\epsilon\), there can be at most \(\lfloor(\Phi_0-\inf\Phi)/\epsilon\rfloor\) such steps. This follows by summing the decreases. Concurrent moves must satisfy the decrease jointly. A workload or capacity change can reset the comparison, and control-plane retries are outside the bound. The result limits repeated accepted moves within the stated regime. It does not bound transient queueing or prove recovery from an overloaded regime.

\paragraph{Worked example: peak backlog versus eventual relief.}
Consider two fluid queues, initially empty, with capacities of 100 records/s each. A four-second burst offers the source 110 records/s, and the destination receives 50 records/s. A movable group accounts for 30 records/s of the source burst. Assume eager migration takes the first two seconds and lowers the source's available foreground service to 85 records/s during the copy. Traffic shifts only after the handoff. For clarity, assume no retries, drops, or further fencing delays and enough destination capacity throughout.

Under fixed placement, the source backlog grows at 10 records/s and peaks at 40 records at the end of the burst. Under eager migration, it grows at 25 records/s during the copy and reaches 50 records at handoff. The shifted placement leaves 80 records/s of load at each endpoint. The source then drains its pending records at 20 records/s until the burst ends, leaving 10 records. Thus the eager policy improves the end-of-burst backlog while raising the peak by 25\%. The saturation-aware rule defers this candidate at burst onset because the foreground demand already exceeds source capacity. With all other decisions held fixed, the resulting backlog equals the fixed-placement trace.

For a burst of length \(H\) and copy duration \(\tau\le H\), the peak comparison for these rates reverses at \(25\tau=10H\). The eager copy has the larger peak if \(\tau>0.4H\). This calculation isolates the feedback risk from lost capacity and uses assumed rates, not system measurements. As utilization approaches one, stochastic queueing can further magnify a service-rate loss~\cite{kingman1961single}. Section~\ref{sec:disturbance-evaluation} tests the distinction between peak harm and later relief by varying burst duration and copy cost separately.

\subsection{Cross-Domain Case Studies}

In the serving case, follow both the request admitted after reclamation and the user who returns to a missing prefix. In the stream case, interrupt a handoff and check which worker resumes from the committed position. In the retrieval case, issue a query while a replica has yet to publish an update. Appendix~C of the supplement gives the event sequences and measurements for each case.

\section{Evaluation and Research Outlook}
\label{sec:evaluation-outlook}

Choose the observation window based on the cost the mechanism can shift. Start at the last comparable state before the action. Continue through transfer, refresh, or eviction, then through the request or update that uses the resulting state. Report useful work as well as queue or memory pressure, so a policy cannot appear to improve by turning work away or making it invisible to the measurement.

\subsection{Evaluation Dimensions for Future Surveys and Systems}\label{sec:disturbance-evaluation}

Measure four quantities together: steady-state throughput and resource use, tail latency and backlog during a disturbance, quality or correctness at the service boundary, and the cost and duration of recovery. Define the service-level objective (SLO) before comparing policies. The relevant quality measure can range from bounded approximation error to retrieval recall or accuracy on earlier tasks.

In Table~\ref{tab:disturbance-evaluation-matrix}, each row pairs the same offered trace with alternative decisions. Align the initial state and the observation window. Follow pending work into recovery and, where relevant, into a later revisit or update. With these controls, an apparent gain in throughput or occupancy can be tested against the work deferred or rejected.

\begin{table}[t]
\small
\centering
\setlength{\tabcolsep}{3pt}
\caption{Proposed disturbance tests. Each measurement spans the disturbance and the work it leaves behind.}
\label{tab:disturbance-evaluation-matrix}
\begin{tabular}{@{}p{0.18\textwidth}p{0.36\textwidth}p{0.40\textwidth}@{}}
\toprule
Scenario & Workload and comparison & What to measure \\
\midrule
Hot-key migration & Vary burst duration and copy cost independently. Compare fixed placement, eager migration, and guarded batches on the same arrivals. & Peak and end-of-burst backlog, late and rejected work, useful completions, migration bytes, reversals, and recovery time. \\
KV reclamation & Fill the cache, then issue requests that revisit the evicted prefixes. Compare against a run with the same arrivals and a retained prefix. & Pages freed during the burst, later reload or prefill time, and time to first token. \\
Index maintenance & Replay the same queries and updates while varying when repair runs and when a version becomes visible. & Recall, query latency, time until a deletion is visible, and pending updates. \\
Sample retention & Shrink memory at the same update in identical task streams. & Actual memory use, examples removed or compressed, and accuracy on earlier tasks after later updates. \\
Shared middleware & Run the same serving, retrieval and update workloads separately and together on a shared link. & Foreground SLO violations, background progress, and unfinished work when the burst ends. \\
\bottomrule
\end{tabular}
\end{table}

The migration and cache tests extend the two observations above: a transfer uses sender resources while live events are arriving~\cite{nasir2019megaphone}, and a remote cache hit still takes time to fetch~\cite{hu2024memserve}. For the migration test in Table~\ref{tab:disturbance-evaluation-matrix}, use the same initial queues, destination capacity, candidate keys, and foreground admission policy. Return the input rate to its pre-burst level and continue measuring until the queues recover or a stated observation limit is reached. Report accepted, deferred and rejected events separately. Compute amplification from peak queues over the same interval and report the absolute peak difference if the fixed-placement peak is zero. Count ownership reversals separately from pauses and retries. Report peak and end-of-burst backlog separately, along with useful completions during the drain. Otherwise a queue that shrinks because requests were dropped can resemble one that drains because service improved.

\subsection{Boundary Conditions and Threats to Transfer}

A transfer claim must name the state object, the correctness or quality constraint, and the measurement horizon. A stream shard has a single writer at each committed epoch, so the key test is recovery on either side of the ownership change. Multiple workers can read a compatible prefix, so the serving test follows a return visit through lookup, reload, or recomputation. An index can meet a version constraint while losing recall under corpus drift. Its evaluation therefore measures both visibility and retrieval quality while updates run.

Use the time scale of the effect to choose the stopping point. A benchmark that covers a copy or reclamation pause may still miss the subsequent rebuild, returning request, or training round. Report which work remains when observation ends, and keep limits on error, freshness, and latency explicit when comparing systems.

\subsection{An Integrated Research Agenda}

\emph{Admitting control work with stale observations.} Collect arrivals, foreground completions, and reservation allocations on the same clock to separate inherited backlog from copy-induced backlog. The unresolved question is whether a conservative demand envelope can be estimated early enough to allow beneficial moves. Vary telemetry age and burst duration independently, then report both violations and useful moves that the guard deferred.

\emph{Coordinating independent data paths.} A scheduling interface can expose lookup, pinning, and movement while leaving the bandwidth decision to several controllers. Each can observe the same free capacity and spend it at the same time. Compare independently admitted copies with shared reservations under synchronized bursts. Measure foreground SLO violations, admitted transfer rate, and starvation of other tenants.

\emph{Estimating the value of state after reclamation.} Cache hits reveal reuse only for objects that survived. Match returning requests to the eviction that affected them, record all reconstruction paths, and measure the latency of later requests. For a replay buffer, trace class coverage and earlier-task accuracy after the budget changes. The unresolved problem is inferring the value of discarding an object when the policy itself changes what can subsequently be observed.

\section{Conclusion}
\label{sec:conclusion}

Across the surveyed systems, state decisions shift work across workers, execution stages, and time. A transferable mechanism must preserve the rules for using the object, account for the resources it shares with foreground work, and follow any cost that returns at reuse. The migration case shows how these requirements yield an admission rule and a backlog bound under stated traffic and service conditions. The remaining challenge is to estimate those conditions under changing demand and to measure the later effects of choices made under a short-term resource constraint.

\clearpage

\bibliographystyle{ACM-Reference-Format}
\bibliography{refs}

\end{document}
